\section{The denominator condition on the original family selection}

% BEGIN X87S1
\paragraph{X87S1. A selected bundle determines a witness over its actual norm and scale.}
Let $\mathcal A$ and $\sigma$ be the built class and its bundle selection.
Use the norm extension and bundle linkage of the common construction.
For each $F\in\mathcal A$, there is one parameter $\kappa$ such that,
for every $h:F.x_1<x$ and $i\in F.I_x$, the following holds.
Put $D=F.D(x,i)$, $O=\mathrm{objAt}(F,x,h,i)$, $c=F.C_{\rm sz}$,
and $b=\sigma(F,x,h,i)$. There are a context $K$ over $\kappa$,
a weighted layer point $q$ in $K$, and a bundle $b_q$ over $q$
with a bundle witness, for which
\[
 (D,O,b)=(q.w.D(q.Z),q.O,b_q),\qquad
 O.L.\mathrm{sc}=K.s,\qquad K.s.x=x .
\]
In particular the norm in the layer object is the norm
$F.\mathrm{ex}(x,h,i).\mathrm{nch}$ in the denominator test.

Indeed, the defining selection condition first gives $\kappa$ for
the entire family. For the chosen $x,h,i$, the assertion that $b$
is built gives $K$, a built point $\beta$ in $K$, its bundle $b'$,
the relation $\kappa.\mathrm{Of}(K)$, and
$(D,O,b)=(\beta.D,\beta.O,b')$. Unpack its linkage. It supplies
$q,b_q$, a bundle witness for these two, and
$(\beta.D,\beta.O,b')=(q.w.D(q.Z),q.O,b_q)$. Compose.
Both equalities live in the sum of pairs $(D,O)$ equipped with a
bundle over that pair. Thus substitution transports the extra data,
their norm component, and the chosen bundle together. In particular
the recursive, division and spectral comparisons in the bundle
witness concern $b$ itself. The scale equality is part of $\beta$,
since $\beta.O$ is the object of its ledger point; its size data
have scale $K.s$. The family size axiom gives
$(F.L(x,i)).\mathrm{sc}.x=x$ under $h$.

The original construction has this defining selection condition:
membership in its class supplies $\kappa$ with the parameter link
and the seven-duty object predicate. Its positivity clause applies
at $F.C_{\rm sz}$ and, at every guarded member, supplies a bundle
which is built and has the joint row property. The selector chooses
a witness of this very existential. Forgetting conjuncts gives the
weaker built condition used above. Thus the full-object equality
requires no nonvanishing assertion. Applying a result that requires
a special range of bases also requires the range argument below.
% END X87S1

% BEGIN X87S2
\paragraph{X87S2. Forgetting the two extra norm conditions fixes the selected object.}
Let $N$ denote the extension with the list of local types. It occurs
in the proof at a layer point. The common bundle construction uses
the extension $N'$ obtained by adjoining first the origin of the
chosen norm and then the prescribed property of the base, frame
and twist. For fixed layer coordinates and norm $a$ its elements
have the form $((n,o),u)$ with $n\in N(\ldots,a,\ldots)$.
Projection to $n$ preserves every index, including $a$ and its
auxiliary data. A linked row established for $N$ therefore applies
with the same bundle and the same selector to the construction with
$N'$.

To see this at the level of witnesses, keep the norm core and
replace only its extension component by this projection. A weighted
point retains its base, frame, twist, columns, layer and full object.
The same substitution at a built core preserves its kernel entry
condition; at a good core it also preserves its goodness condition.
In the bundle witness keep the twisted input, the recursive and kernel
comparisons, the spectral and point equalities, the summand and
division pins, the side and the column pins. The recursive comparison
over the weaker extension has the same host, columns, forms,
characters, Galois data, kernel and coordinate fields, copied from
the original comparison. All these fields use the norm core retained
by projection. This gives the linkage at $N$ for the very triple
over which linkage at $N'$ was given.

For an additional slot $R$ allowed over built cores for $N$, use
the identical $R$ for $N'$: each built core for $N'$ projects to
one for $N$. A good core with $R$ for $N'$ projects to a good core
with $R$ for $N$. The eventual assertion at such cores now follows
by applying its old implication and the preceding comparison of
linkages. A selected bundle realizing the stronger linkage realizes
the weaker linkage at the same $K,\beta$; the projection on bundles
is the identity. Consequently the same row predicate, the same
$\sigma$, and its previously obtained conclusion work for $N'$.
This is the norm change used before the denominator row joins the
other rows.
% END X87S2

% BEGIN X87S3
\paragraph{X87S3. How an eventual layer theorem becomes a condition on the chosen bundle.}
Write $P(D,a)$ for a property depending on the datum and its norm,
and assume its eventual layer assertion on the source's default
extension, row extension and column pin. Write $R_P$ for the
condition at a fixed base, frame and twist that this assertion holds
for all their qualifying layer data. For the extension with list,
the built and good slots are denoted by $Q_b,Q_g$. There is a
permitted additional slot and an eventual row on it whose object
predicate is precisely $E_P(D,O,c,b)=P(D,O.\mathrm{ex}.\mathrm{nch})$.

Here are the two quantifiers involved in this passage. For a context
$K$, let $G_P(K)$ assert that on sufficiently general bases, frames
and twists over $K$ satisfying the base condition, every core in
$Q_b$ satisfies $R_P$. This uses the source's large-base quantifier,
including its admissible avoided sets. By the layer assertion,
$G_P(K)$ holds eventually among the good contexts. The additional
slot used before choosing a context is
\[
       G_P(K)\ \Longrightarrow\ R_P(K,X,f,T,c_0).
\]
It is allowed in the unconditioned large-base quantifier: if $G_P(K)$
holds, use the admissible range that it gives; otherwise the
implication holds at every base. On the eventual good contexts
where $G_P(K)$ holds, intersect its range with the range used
for the row. Thus a good core in the additional slot supplies the
actual $R_P$ there.

For a built point $\beta$ in this intersection, take a building
witness from its good core and equip it with that core's goodness.
The ledger point of $\beta$ uses the division lattice
$n'(\beta.p.L)$ and height $\beta.p.\mathrm{ht}$. Move the witness
to these using the division transport of X86B2. The division
identity for the datum and the ledger identity for the object
give a weighted point $q$ with
$(\beta.D,\beta.O)=(q.w.D(q.Z),q.O)$. The list extension projects to
the law extension and then to the default extension on the same norm.
The witness contains the row extension and the column pin.
Evaluating $R_P$ on these inputs gives $P(q.w.D(q.Z),q.O.\mathrm{ex}.\mathrm{nch})$.
Transport along the pair equality gives $E_P$ for $\beta$ and every
linked bundle at it. This proves the asserted eventual row. Its
object predicate records the result after the good-context
implication has been discharged. No bound on the bases of all
individual witnesses in X87S1 is needed.
% END X87S3

% BEGIN X87S4
\paragraph{X87S4. The denominator row reaches the class used in the final window argument.}
Take $P$ in X87S3 to be the eighteen original preconditions together
with the four ties. Retain the hypotheses of the family implication
in A85D3. Suppose the layer inequality of F86D1 holds on the
eventual range in which the other member assertions of F86D2 hold.
Then the denominator contributes the required row to the original
joint class construction. In particular, on the common class
produced by that construction the exact implication is
\[
    \mathrm{Closed}\bigl(\mathrm{Des}(\mathrm{WithChN}(\mathcal A),3),
                         \,\mathrm{LinkN}\bigr).
\]
This conclusion uses the usual common row and linked-bundle
existence inputs when producing the common family. No further
identification of selectors is a hypothesis of this implication.

Indeed F86D1--2 give the missing two fields from the stated
inequality. Intersect the finitely many eventual ranges for the
fields and the four ties, retaining their order of parameters.
Apply X87S3. As a row valid for selections, take its object
predicate to be $E_P$ itself. If $\sigma$ realizes it, its defining
condition at each $x>F.x_1$ and every $i$ supplies
$P(F.D(x,i),F.\mathrm{ex}(x,h,i).\mathrm{nch})$.
Membership in $D_3(3,\mathcal A)=\mathrm{Des}(\mathrm{WithChN}(\mathcal A),3)$
includes both membership in $\mathcal A$ and the twist condition.
Apply X86B1 to that very family with $b=F.x_1$, giving its
denominator bound with the original cohomological finiteness and
cyclotomic summability guards. Equivalently one can retain the
source's additional row at $2,3$: intersect the permitted slots,
project to $E_P$, and take this second row's outer exponent and
transport parameters before $F$ and $C_H$ are given, as in X86B2.

It remains to locate this row in the joint selection. First transport
it by X87S2 to the norm extension of the common bundle. The source
places it in the entry $\mathrm{nRest}$ of the record of rows.
For two such rows with permitted slots $R_1,R_2$, take
$R_1\wedge R_2$. Its large-base assertion follows by intersecting
the two admissible ranges. Strengthen both row implications by
the relevant projections and use the conjunction of their
object predicates. Iterate over the finite row record, including
$\mathrm{nRest}$, the count row with its already chosen constants,
and the row for $c_T(\varepsilon,y)=c_{T,\min}(y)$. Intersect once
more with the permitted slot for existence of a linked bundle.
Write $R$ for the resulting extra slot, $L$ for the common linkage,
and $E$ for the conjoined object predicate. On the common eventual
range, first choose a linked bundle at the given point and size
constant; the universal row implications apply to this choice.
This supplies the object property for that size and hence for each
size tested by the class definition.

Now unpack the source's single call to class selection. Its class
and its selector are, respectively,
\[
    \mathcal A=\mathrm{cls}(Q_g'\wedge R,L,E),\qquad
    \sigma=\mathrm{pick}(Q_g'\wedge R,L,E),
\]
where $Q_g'$ is the good slot for the stronger extension.
The seven-duty predicate defining $\mathcal A$ supplies at each
guarded member a bundle which is built and satisfies $E$, as in
X87S1. Thus every projection of $E$, including the denominator's
predicate, concerns this same choice. This proves the conclusion
without using $\mathrm{nRest}$ as an assumption: its member proof
was supplied before forming the conjunction. Existence of a
target family uses the separate recursive and geometric inputs of
the common design. In the final window argument the host restriction
replaces $\mathcal A$ by $\mathcal A\cap\mathcal E$ for the
design-existence input. More explicitly put $\mathcal C'=\mathrm{WithChN}(\mathcal A\cap\mathcal E)$. A membership proof for $\mathcal C'$ consists of membership in the intersection and the same auxiliary character conditions. Project the intersection to $\mathcal A$ and keep the auxiliary conditions. This gives $\mathcal C'\subseteq\mathrm{WithChN}(\mathcal A)$ and, on adding the same smoothness, twist, chosen exponent, tuple-size and discriminant-size conditions, the corresponding inclusion of their design classes. The selector on the intersection is $\sigma(F,h.1)$ for the original membership proof $h$ in the intersection. Its value at each guarded member is therefore the original bundle. The denominator input still takes
$\mathcal C=\mathrm{WithChN}(\mathcal A)$. On a target family the
inclusion simply projects $F\in\mathcal A\cap\mathcal E$ to
$F\in\mathcal A$, preserving the family, its norms and its
$C_H$. Thus the class at which the denominator estimate is used
is identified by the original construction. Its outstanding
member input is the reverse multiplicity inequality with the
uniform layer quantifiers stated above.
% END X87S4
